\section*{Training 2. Two shocks, one average · 35 points}
Two wealth lotteries have equally likely outcomes:
\[
 A:\ (36,100),\qquad B:\ (4,132).
\]
Use Bernoulli indices $b(w)=\sqrt w$ and $g(w)=\log w$ on $w>0$.
The certainty equivalent solves $b(CE)=\mathbb E[b(W)]$; the risk premium
is $\mathbb E[W]-CE$.
\begin{enumerate}[label=(\alph*)]
\item \textbf{6 points.} Calculate both means. Compare their cumulative
distribution functions and determine whether either lottery first-order
stochastically dominates the other.
\item \textbf{8 points.} Under b, calculate expected utility, the certainty
equivalent and the risk premium for each lottery.
\item \textbf{7 points.} Find A's certainty equivalent under g. Compare A
with sure wealth 62 under each index. Explain why an increasing change in
the Bernoulli index can change lottery rankings.
\item \textbf{6 points.} Calculate absolute and relative risk aversion for
b and g. Which index represents greater absolute risk aversion at each wealth?
\item \textbf{8 points.} Use Jensen's inequality to prove that a consumer
with an increasing concave Bernoulli index has a nonnegative risk premium.
When does strict concavity imply a strictly positive premium?
\end{enumerate}
\paragraph{Further practice, ungraded.} Couple the lotteries as follows:
\[
 \Pr(B=4\mid A=36)=3/4,\qquad \Pr(B=4\mid A=100)=1/4.
\]
Check B's unconditional distribution and show $\mathbb E[B\mid A]=A$.
Apply conditional Jensen to show that every increasing concave Bernoulli
index weakly prefers A. This establishes second-order stochastic dominance.

\clearpage
\section*{Training 3. The umbrella with a loading · 35 points}
A household has wealth 100 and faces a loss of 84 with probability $1/4$.
Its Bernoulli utility is $b(w)=\log w$. Coverage $I\in[0,84]$ pays I only
if the loss occurs. The premium is $cI$, paid in either state.
\begin{enumerate}[label=(\alph*)]
\item \textbf{6 points.} Write both final wealth levels and expected utility.
Explain existence and uniqueness of optimal coverage for $c=1/4$ and $c=1/2$.
\item \textbf{7 points.} At the actuarially fair price $c=1/4$, find optimal
coverage. Give a Jensen argument, not only a stationary-point calculation.
\item \textbf{10 points.} At $c=1/2$, find optimal coverage and both final
wealth levels. Check that your candidate lies within the allowed interval
and is globally optimal.
\item \textbf{12 points.} Keep the loss 84, probability $1/4$ and $c=1/2$,
but replace initial wealth by $W>84$. Derive optimal coverage as a function
of W, including the no-insurance region. Relate the result to absolute risk
aversion for log utility.
\end{enumerate}
